% Part: normal-modal-logic
% Chapter: sequent-calculus
% Section: introduction

\documentclass[../../../include/open-logic-section]{subfiles}

\begin{document}

\olfileid{nml}{seq}{int}

\olsection{పరిచయం}

ప్రతిజ్ఞావాక్య తర్కపు సీక్వెంట్ కలనశాస్త్రాన్ని
\iftag{prvBox}{$\Box$\iftag{prvDiamond}{ మరియు~}{}}{}%
\iftag{prvDiamond}{$\Diamond$}{} కోసం అదనపు నియమాలతో
విస్తరించవచ్చు. ఉదాహరణకు, $\Log{K}$ కోసం
\Log{LK}తో పాటు ఇవి ఉంటాయి:
\[
  \iftag{prvBox}{\iftag{prvDiamond}{% <> and [] primitive
    \Axiom$\Gamma \fCenter \Delta, !A$
    \RightLabel{$\Box$}
    \UnaryInf$\Box\Gamma \fCenter \Diamond\Delta, \Box!A$
    \DisplayProof
    \qquad
    \Axiom$!A, \Gamma \fCenter \Delta$
    \RightLabel{$\Diamond$}
    \UnaryInf$\Diamond!A, \Box\Gamma \fCenter \Diamond\Delta$
    \DisplayProof
}{% only Box primitive
\Axiom$\Gamma \fCenter !A$
\RightLabel{$\Box$}
\UnaryInf$\Box\Gamma \fCenter \Box!A$
\DisplayProof
}}{% only <> primitive
  \Axiom$!A \fCenter \Delta$
  \RightLabel{$\Diamond$}
  \UnaryInf$\Diamond!A \fCenter \Diamond\Delta$
  \DisplayProof
}
\]
$\Log{K}$ విస్తరణలకు ఇంకా అదనపు నియమాలు చేర్చాలి.

ప్రతి మోడల్ తర్కానికీ ఇలాంటి సీక్వెంట్ కలనశాస్త్రం
ఉండదు. అర్థవిజ్ఞాన పరంగా సరళమైన $\Log{S5}$ను
ప్రాప్యత సంబంధాలను అసలు ఉపయోగించకుండానే
నిర్వచించవచ్చు. అయినా, $\Log{LK}$ నుంచి ఏర్పడి,
\Cut{} నియమం లేకుండానే సంపూర్ణంగా ఉండే సీక్వెంట్
కలనశాస్త్రం దానికి ఉన్నట్టు తెలియదు. అయితే దానికి
కట్-రహిత, సంపూర్ణ \emph{హైపర్ సీక్వెంట్}
కలనశాస్త్రం ఉంది.

\end{document}
